%!TEX root = ../main.tex
% Chapitre 7 : Transformers
\chapter{Transformers}
\label{ch:transformers}

\paragraph{Objectifs du chapitre}
\begin{itemize}
\item Comprendre l'architecture Transformer et le m\'ecanisme d'attention
\item D\'eriver les formules de l'attention multi-t\^ete et de l'encodage positionnel
\item Distinguer les variantes encodeur, d\'ecodeur et encodeur-d\'ecodeur
\item Impl\'ementer un Transformer complet en PyTorch
\item Conna\^itre les mod\`eles fondateurs : BERT, GPT, ViT
\end{itemize}

\section{<<~Attention Is All You Need~>> --- Le changement de paradigme}
\label{sec:transformers-intro}

En 2017, Vaswani et al.\ publient <<~Attention Is All You Need~>>\index{Attention Is All You Need}, un article qui transforme radicalement le paysage du traitement automatique des langues et, plus largement, de l'apprentissage profond\index{Transformer}. L'id\'ee centrale est de remplacer enti\`erement les r\'ecurrences (RNN, LSTM, GRU --- cf.\ chapitre pr\'ec\'edent) par un m\'ecanisme d'\emph{attention}\index{attention} pur, permettant un parall\'elisme massif et une mod\'elisation efficace des d\'ependances \`a longue port\'ee.

\begin{remarque}
Avant les Transformers, les mod\`eles s\'equence-\`a-s\'equence\index{sequence-to-sequence} reposaient sur des encodeurs-d\'ecodeurs RNN avec attention (Bahdanau, 2014). Le Transformer \'elimine toute r\'ecurrence, ne conservant que l'attention.
\end{remarque}

\section{M\'ecanisme d'attention}
\label{sec:attention-mechanism}

\subsection{Attention Scaled Dot-Product}

Soit une s\'equence d'entr\'ee de longueur $n$ avec des vecteurs de dimension $d_k$. On d\'efinit trois projections lin\'eaires :
\begin{itemize}
    \item \textbf{Query} (requ\^ete) : $Q = XW^Q$, avec $W^Q \in \R^{d_{\text{model}} \times d_k}$
    \item \textbf{Key} (cl\'e) : $K = XW^K$, avec $W^K \in \R^{d_{\text{model}} \times d_k}$
    \item \textbf{Value} (valeur) : $V = XW^V$, avec $W^V \in \R^{d_{\text{model}} \times d_v}$
\end{itemize}

\begin{mathbox}[title=\textbf{Attention Scaled Dot-Product}]
\begin{equation}
\text{Attention}(Q, K, V) = \softmax\!\left(\frac{QK^\top}{\sqrt{d_k}}\right) V
\label{eq:scaled-attention-transformer}
\end{equation}
\end{mathbox}

\begin{intuition}
Le facteur $\frac{1}{\sqrt{d_k}}$ emp\^eche les produits scalaires de devenir trop grands en dimension \'elev\'ee, ce qui pousserait le $\softmax$ dans des r\'egions de gradient quasi nul.

En effet, si les composantes de $Q$ et $K$ sont i.i.d.\ de moyenne 0 et variance 1, alors $\E[q_i^\top k_j] = 0$ et $\Var(q_i^\top k_j) = d_k$. La division par $\sqrt{d_k}$ ram\`ene la variance \`a 1.
\end{intuition}

\subsection{Attention multi-t\^ete}
\label{sec:multi-head-attention}

Plut\^ot qu'une seule fonction d'attention avec $d_{\text{model}}$ dimensions, on utilise $h$ <<~t\^etes~>>\index{multi-head attention} parall\`eles :

\begin{formules}[title=\textbf{Attention multi-t\^ete}]
\begin{align}
\text{MultiHead}(Q, K, V) &= \text{Concat}(\text{head}_1, \ldots, \text{head}_h)\, W^O \\
\text{o\`u} \quad \text{head}_i &= \text{Attention}(QW_i^Q,\; KW_i^K,\; VW_i^V)
\end{align}
avec $W_i^Q \in \R^{d_{\text{model}} \times d_k}$, $W_i^K \in \R^{d_{\text{model}} \times d_k}$, $W_i^V \in \R^{d_{\text{model}} \times d_v}$, $W^O \in \R^{hd_v \times d_{\text{model}}}$, et typiquement $d_k = d_v = d_{\text{model}} / h$.
\end{formules}

\begin{remarque}
Chaque t\^ete peut apprendre \`a se concentrer sur un type diff\'erent de relation : certaines t\^etes captent la syntaxe, d'autres la s\'emantique, d'autres encore les relations positionnelles.
\end{remarque}

\section{Encodage positionnel}
\label{sec:positional-encoding}

Le Transformer n'ayant pas de r\'ecurrence, il n'a aucune notion intrins\`eque de l'ordre des tokens. On ajoute donc un \emph{encodage positionnel}\index{encodage positionnel} aux embeddings d'entr\'ee.

\subsection{D\'erivation de la formule sinuso\"idale}

\begin{definition}[Encodage positionnel sinuso\"idal]
Pour une position $\text{pos}$ et une dimension $i$ :
\begin{align}
PE_{(\text{pos}, 2i)} &= \sin\!\left(\frac{\text{pos}}{10000^{2i/d_{\text{model}}}}\right) \label{eq:pe-sin}\\
PE_{(\text{pos}, 2i+1)} &= \cos\!\left(\frac{\text{pos}}{10000^{2i/d_{\text{model}}}}\right) \label{eq:pe-cos}
\end{align}
\end{definition}

\begin{theoreme}[Propri\'et\'e de position relative]
Pour tout d\'ecalage fixe $k$, il existe une transformation lin\'eaire $M_k$ ind\'ependante de $\text{pos}$ telle que :
\begin{equation}
PE_{\text{pos}+k} = M_k \cdot PE_{\text{pos}}
\end{equation}
\end{theoreme}

\begin{proof}
Consid\'erons la paire de dimensions $(2i, 2i+1)$ et posons $\omega_i = 1/10000^{2i/d_{\text{model}}}$. On a :
\begin{align}
\begin{pmatrix} \sin(\omega_i(\text{pos}+k)) \\ \cos(\omega_i(\text{pos}+k)) \end{pmatrix}
&= \begin{pmatrix} \cos(\omega_i k) & \sin(\omega_i k) \\ -\sin(\omega_i k) & \cos(\omega_i k) \end{pmatrix}
\begin{pmatrix} \sin(\omega_i \cdot \text{pos}) \\ \cos(\omega_i \cdot \text{pos}) \end{pmatrix}
\end{align}
par les formules d'addition trigonom\'etriques. La matrice de rotation ne d\'epend que de $k$, pas de $\text{pos}$, ce qui permet au mod\`ele d'apprendre \`a exploiter les positions relatives.
\end{proof}

\section{Architecture compl\`ete du Transformer}
\label{sec:transformer-architecture}

\subsection{Bloc encodeur}
\label{sec:encoder-block}

Chaque bloc encodeur est compos\'e de :
\begin{enumerate}
    \item \textbf{Multi-Head Self-Attention} (attention sur soi-m\^eme)
    \item \textbf{Add \& Norm} (connexion r\'esiduelle + normalisation de couche)
    \item \textbf{Feed-Forward Network} (FFN) \`a deux couches :
    \begin{equation}
    \text{FFN}(x) = \relu(xW_1 + b_1)W_2 + b_2
    \end{equation}
    avec $W_1 \in \R^{d_{\text{model}} \times d_{ff}}$, $W_2 \in \R^{d_{ff} \times d_{\text{model}}}$, et typiquement $d_{ff} = 4 \cdot d_{\text{model}}$.
    \item \textbf{Add \& Norm} (seconde connexion r\'esiduelle + normalisation)
\end{enumerate}

\subsection{Bloc d\'ecodeur}
\label{sec:decoder-block}

Le d\'ecodeur ajoute une couche suppl\'ementaire :
\begin{enumerate}
    \item \textbf{Masked Multi-Head Self-Attention}\index{masked attention} : un masque causal emp\^eche chaque position d'<<~observer~>> les positions futures.
    \item \textbf{Add \& Norm}
    \item \textbf{Multi-Head Cross-Attention}\index{cross-attention} : les queries viennent du d\'ecodeur, les keys et values viennent de la sortie de l'encodeur.
    \item \textbf{Add \& Norm}
    \item \textbf{FFN}
    \item \textbf{Add \& Norm}
\end{enumerate}

\subsection{Pre-Norm vs Post-Norm}
\label{sec:pre-post-norm}

\begin{definition}[Post-Norm (original)]
\begin{equation}
x_{l+1} = \text{LayerNorm}(x_l + \text{SubLayer}(x_l))
\label{eq:post-norm}
\end{equation}
\end{definition}

\begin{definition}[Pre-Norm]
\begin{equation}
x_{l+1} = x_l + \text{SubLayer}(\text{LayerNorm}(x_l))
\label{eq:pre-norm}
\end{equation}
\end{definition}

\begin{proposition}[Stabilit\'e du gradient avec Pre-Norm]
Dans le cas Pre-Norm, le gradient traverse directement la connexion r\'esiduelle sans passer par la normalisation :
\begin{equation}
\frac{\partial x_L}{\partial x_l} = I + \sum_{k=l}^{L-1} \frac{\partial \text{SubLayer}_k(\text{LayerNorm}(x_k))}{\partial x_l}
\end{equation}
Le terme identit\'e $I$ garantit un flux de gradient non nul, ce qui stabilise l'entra\^inement pour les architectures profondes ($L > 12$).
\end{proposition}

\begin{attention}
Le Transformer original utilise Post-Norm, mais la majorit\'e des impl\'ementations modernes (GPT-2, GPT-3, LLaMA) utilisent Pre-Norm pour sa stabilit\'e sup\'erieure. Un \emph{warm-up} du taux d'apprentissage est g\'en\'eralement n\'ecessaire avec Post-Norm.
\end{attention}

\subsection{Diagramme TikZ de l'architecture compl\`ete}

\begin{figure}[htbp]
\centering
\begin{tikzpicture}[
    block/.style={rectangle, draw, fill=blue!8, minimum width=3.2cm, minimum height=0.8cm, align=center, rounded corners=2pt},
    attnblock/.style={rectangle, draw, fill=orange!15, minimum width=3.2cm, minimum height=0.8cm, align=center, rounded corners=2pt},
    ffnblock/.style={rectangle, draw, fill=green!12, minimum width=3.2cm, minimum height=0.8cm, align=center, rounded corners=2pt},
    normblock/.style={rectangle, draw, fill=gray!15, minimum width=3.2cm, minimum height=0.6cm, align=center, rounded corners=2pt},
    io/.style={rectangle, draw, fill=yellow!15, minimum width=3.2cm, minimum height=0.7cm, align=center, rounded corners=2pt},
    arrow/.style={-{Stealth}, thick},
    dashed arrow/.style={-{Stealth}, thick, dashed},
    brace/.style={decorate, decoration={brace, amplitude=6pt, raise=4pt}},
]

% Encodeur
\node[io] (emb-enc) at (0, 0) {Input Embedding\\+ Pos.\ Encoding};
\node[attnblock] (sa-enc) at (0, 1.5) {Multi-Head\\Self-Attention};
\node[normblock] (an1-enc) at (0, 2.6) {Add \& Norm};
\node[ffnblock] (ffn-enc) at (0, 3.6) {Feed-Forward};
\node[normblock] (an2-enc) at (0, 4.7) {Add \& Norm};

\draw[arrow] (emb-enc) -- (sa-enc);
\draw[arrow] (sa-enc) -- (an1-enc);
\draw[arrow] (an1-enc) -- (ffn-enc);
\draw[arrow] (ffn-enc) -- (an2-enc);

% Accolade encodeur
\draw[brace] (2.0, 1.1) -- (2.0, 5.1) node[midway, right=10pt, align=center] {$\times N$};

% D\'ecodeur
\node[io] (emb-dec) at (7, 0) {Output Embedding\\+ Pos.\ Encoding};
\node[attnblock] (msa-dec) at (7, 1.5) {Masked Multi-Head\\Self-Attention};
\node[normblock] (an1-dec) at (7, 2.6) {Add \& Norm};
\node[attnblock] (ca-dec) at (7, 3.6) {Multi-Head\\Cross-Attention};
\node[normblock] (an2-dec) at (7, 4.7) {Add \& Norm};
\node[ffnblock] (ffn-dec) at (7, 5.7) {Feed-Forward};
\node[normblock] (an3-dec) at (7, 6.8) {Add \& Norm};
\node[io] (out) at (7, 7.9) {Linear + Softmax};

\draw[arrow] (emb-dec) -- (msa-dec);
\draw[arrow] (msa-dec) -- (an1-dec);
\draw[arrow] (an1-dec) -- (ca-dec);
\draw[arrow] (ca-dec) -- (an2-dec);
\draw[arrow] (an2-dec) -- (ffn-dec);
\draw[arrow] (ffn-dec) -- (an3-dec);
\draw[arrow] (an3-dec) -- (out);

% Accolade d\'ecodeur
\draw[brace] (9.0, 1.1) -- (9.0, 7.2) node[midway, right=10pt, align=center] {$\times N$};

% Cross-attention arrow
\draw[dashed arrow] (an2-enc.east) -- ++(1.0, 0) |- (ca-dec.west)
    node[midway, above, font=\footnotesize] {$K, V$};

% Labels
\node[above=0.2cm, font=\bfseries] at (0, 5.1) {Encodeur};
\node[above=0.2cm, font=\bfseries] at (7, 8.3) {D\'ecodeur};

\end{tikzpicture}
\caption{Architecture compl\`ete du Transformer avec encodeur ($\times N$ blocs) et d\'ecodeur ($\times N$ blocs). Les fl\`eches en pointill\'e repr\'esentent le flux d'information de l'encodeur vers le m\'ecanisme de cross-attention du d\'ecodeur.}
\label{fig:transformer-architecture}
\end{figure}

\section{Vision Transformer (ViT)}
\label{sec:vit}

Le \emph{Vision Transformer}\index{Vision Transformer}\index{ViT} (Dosovitskiy et al., 2020) adapte l'architecture Transformer au domaine de la vision par ordinateur.

\begin{definition}[Patch Embedding]
Une image $x \in \R^{H \times W \times C}$ est d\'ecoup\'ee en $N = HW/P^2$ patches de taille $P \times P$, chacun aplati et projet\'e lin\'eairement :
\begin{equation}
z_0 = [\texttt{[CLS]};\; x_1^p E;\; x_2^p E;\; \ldots;\; x_N^p E] + E_{\text{pos}}
\end{equation}
o\`u $E \in \R^{P^2 C \times d_{\text{model}}}$ est la matrice de projection et $E_{\text{pos}} \in \R^{(N+1) \times d_{\text{model}}}$ l'encodage positionnel appris.
\end{definition}

\begin{remarque}
Le token sp\'ecial \texttt{[CLS]}\index{CLS token} joue le r\^ole d'agr\'egateur : sa repr\'esentation finale est utilis\'ee pour la classification.
\end{remarque}

\section{BERT : Mod\'elisation de langage masqu\'ee}
\label{sec:bert}

\textbf{BERT}\index{BERT} (Bidirectional Encoder Representations from Transformers, Devlin et al., 2019) utilise uniquement l'\emph{encodeur} du Transformer.

\begin{definition}[Objectif MLM\index{Masked Language Modeling}]
On masque al\'eatoirement 15\% des tokens d'une s\'equence. Le mod\`ele doit pr\'edire les tokens masqu\'es :
\begin{equation}
\mathcal{L}_{\text{MLM}} = -\E\!\left[\sum_{i \in \mathcal{M}} \log p(x_i \mid x_{\setminus \mathcal{M}};\, \theta)\right]
\end{equation}
o\`u $\mathcal{M}$ est l'ensemble des positions masqu\'ees.
\end{definition}

\section{GPT : Mod\'elisation de langage autorégressive}
\label{sec:gpt}

\textbf{GPT}\index{GPT} (Generative Pre-trained Transformer, Radford et al., 2018) utilise uniquement le \emph{d\'ecodeur} avec un masque causal\index{masque causal}.

\begin{definition}[Objectif autoregressif]
Le mod\`ele pr\'edit chaque token conditionnellement aux tokens pr\'ec\'edents :
\begin{equation}
\mathcal{L}_{\text{AR}} = -\sum_{t=1}^{T} \log p(x_t \mid x_1, \ldots, x_{t-1};\, \theta)
\end{equation}
\end{definition}

Le masque causal est r\'ealis\'e en ajoutant $-\infty$ aux positions futures dans la matrice d'attention avant le $\softmax$ :
\begin{equation}
\text{mask}_{ij} = \begin{cases} 0 & \text{si } j \leq i \\ -\infty & \text{si } j > i \end{cases}
\end{equation}

\section{Impl\'ementation PyTorch d'un Transformer}
\label{sec:transformer-pytorch}

\begin{codeblock}[title=\textbf{Transformer complet pour la classification de s\'equences}]
\begin{minted}{python}
import torch
import torch.nn as nn
import math

class PositionalEncoding(nn.Module):
    """Encodage positionnel sinusoidal."""
    def __init__(self, d_model, max_len=512, dropout=0.1):
        super().__init__()
        self.dropout = nn.Dropout(p=dropout)
        pe = torch.zeros(max_len, d_model)
        position = torch.arange(0, max_len).unsqueeze(1).float()
        div_term = torch.exp(
            torch.arange(0, d_model, 2).float()
            * (-math.log(10000.0) / d_model)
        )
        pe[:, 0::2] = torch.sin(position * div_term)
        pe[:, 1::2] = torch.cos(position * div_term)
        pe = pe.unsqueeze(0)  # (1, max_len, d_model)
        self.register_buffer('pe', pe)

    def forward(self, x):
        # x: (batch, seq_len, d_model)
        x = x + self.pe[:, :x.size(1)]
        return self.dropout(x)


class MultiHeadAttention(nn.Module):
    """Attention multi-tete from scratch."""
    def __init__(self, d_model, n_heads):
        super().__init__()
        assert d_model % n_heads == 0
        self.d_k = d_model // n_heads
        self.n_heads = n_heads
        self.W_q = nn.Linear(d_model, d_model)
        self.W_k = nn.Linear(d_model, d_model)
        self.W_v = nn.Linear(d_model, d_model)
        self.W_o = nn.Linear(d_model, d_model)

    def forward(self, query, key, value, mask=None):
        B, T, D = query.shape
        # Projections et reshape : (B, n_heads, T, d_k)
        Q = self.W_q(query).view(B, T, self.n_heads, self.d_k).transpose(1, 2)
        K = self.W_k(key).view(B, -1, self.n_heads, self.d_k).transpose(1, 2)
        V = self.W_v(value).view(B, -1, self.n_heads, self.d_k).transpose(1, 2)

        # Scaled dot-product attention
        scores = torch.matmul(Q, K.transpose(-2, -1)) / math.sqrt(self.d_k)
        if mask is not None:
            scores = scores.masked_fill(mask == 0, float('-inf'))
        attn = torch.softmax(scores, dim=-1)
        out = torch.matmul(attn, V)

        # Concatenation et projection finale
        out = out.transpose(1, 2).contiguous().view(B, T, D)
        return self.W_o(out)


class TransformerEncoderBlock(nn.Module):
    """Bloc encodeur Pre-Norm."""
    def __init__(self, d_model, n_heads, d_ff, dropout=0.1):
        super().__init__()
        self.attn = MultiHeadAttention(d_model, n_heads)
        self.ffn = nn.Sequential(
            nn.Linear(d_model, d_ff),
            nn.GELU(),
            nn.Dropout(dropout),
            nn.Linear(d_ff, d_model),
            nn.Dropout(dropout),
        )
        self.ln1 = nn.LayerNorm(d_model)
        self.ln2 = nn.LayerNorm(d_model)
        self.drop = nn.Dropout(dropout)

    def forward(self, x, mask=None):
        # Pre-Norm
        x_norm = self.ln1(x)
        x = x + self.drop(self.attn(x_norm, x_norm, x_norm, mask))
        x = x + self.ffn(self.ln2(x))
        return x


class TransformerClassifier(nn.Module):
    """Transformer encodeur pour la classification de sequences."""
    def __init__(self, vocab_size, d_model=128, n_heads=4,
                 d_ff=512, n_layers=4, n_classes=2,
                 max_len=256, dropout=0.1):
        super().__init__()
        self.embedding = nn.Embedding(vocab_size, d_model)
        self.pos_enc = PositionalEncoding(d_model, max_len, dropout)
        self.layers = nn.ModuleList([
            TransformerEncoderBlock(d_model, n_heads, d_ff, dropout)
            for _ in range(n_layers)
        ])
        self.ln_final = nn.LayerNorm(d_model)
        self.classifier = nn.Linear(d_model, n_classes)

    def forward(self, x, mask=None):
        # x: (batch, seq_len) indices entiers
        x = self.embedding(x) * math.sqrt(self.embedding.embedding_dim)
        x = self.pos_enc(x)
        for layer in self.layers:
            x = layer(x, mask)
        x = self.ln_final(x)
        # Moyenne sur la sequence (alternative au token [CLS])
        x = x.mean(dim=1)
        return self.classifier(x)


# --- Entrainement ---
model = TransformerClassifier(
    vocab_size=10000, d_model=128, n_heads=4,
    d_ff=512, n_layers=4, n_classes=2
)
optimizer = torch.optim.AdamW(model.parameters(), lr=1e-4, weight_decay=0.01)
criterion = nn.CrossEntropyLoss()

# Exemple avec donnees synthetiques
batch_x = torch.randint(0, 10000, (32, 64))  # (batch=32, seq_len=64)
batch_y = torch.randint(0, 2, (32,))

logits = model(batch_x)
loss = criterion(logits, batch_y)
loss.backward()
optimizer.step()
print(f"Loss: {loss.item():.4f}, Logits shape: {logits.shape}")
\end{minted}
\end{codeblock}

\begin{codeoutput}
\begin{verbatim}
Loss: 0.6847, Logits shape: torch.Size([32, 2])
\end{verbatim}
\end{codeoutput}

\begin{bonnepratique}
Utilisez toujours \texttt{AdamW}\index{AdamW} avec un \emph{weight decay} pour l'entra\^inement des Transformers. Un \emph{learning rate schedule} avec warm-up lin\'eaire suivi d'un d\'eclin cosinus est standard.
\end{bonnepratique}

\section{\'Etude d'ablation}
\label{sec:transformer-ablation}

\begin{table}[htbp]
\centering
\caption{\'Etude d'ablation sur un Transformer pour la classification de texte (IMDB). Pr\'ecision sur l'ensemble de test (\%).}
\label{tab:transformer-ablation}
\begin{tabular}{lcccc}
\toprule
\textbf{Configuration} & \textbf{Couches} & \textbf{T\^etes} & \textbf{$d_{ff}$} & \textbf{Pr\'ecision (\%)} \\
\midrule
Base (Post-Norm) & 4 & 4 & 512 & 86.3 \\
Base (Pre-Norm) & 4 & 4 & 512 & 87.1 \\
\midrule
1 couche & 1 & 4 & 512 & 82.5 \\
2 couches & 2 & 4 & 512 & 85.4 \\
6 couches & 6 & 4 & 512 & 87.4 \\
8 couches & 8 & 4 & 512 & 87.2 \\
\midrule
1 t\^ete & 4 & 1 & 512 & 84.8 \\
2 t\^etes & 4 & 2 & 512 & 86.0 \\
8 t\^etes & 4 & 8 & 512 & 87.3 \\
\midrule
$d_{ff}=128$ & 4 & 4 & 128 & 84.1 \\
$d_{ff}=256$ & 4 & 4 & 256 & 85.9 \\
$d_{ff}=1024$ & 4 & 4 & 1024 & 87.5 \\
\bottomrule
\end{tabular}
\end{table}

\begin{remarque}
Les r\'esultats montrent que : (1) Pre-Norm est syst\'ematiquement sup\'erieur \`a Post-Norm sans warm-up fin ; (2) augmenter le nombre de couches au-del\`a de 6 apporte peu pour ce jeu de donn\'ees de taille modeste ; (3) la dimension du FFN a un impact significatif ; (4) une seule t\^ete d'attention est nettement insuffisante.
\end{remarque}

\section{\'Etat de l'art et connexions}
\label{sec:transformers-sota}

\begin{sota}[title=\textbf{Grands mod\`eles de langage (LLMs)}]\index{LLM}\index{large language model}
Les Transformers sont \`a la base de la r\'evolution des LLMs :

\begin{itemize}
    \item \textbf{Lois d'\'echelle (Scaling Laws)}\index{scaling laws} : Kaplan et al.\ (2020) montrent que la performance d'un LLM suit une loi de puissance en fonction du nombre de param\`etres $N$, de la taille du jeu de donn\'ees $D$ et du budget de calcul $C$ :
    \begin{equation}
    L(N) \approx \left(\frac{N_c}{N}\right)^{\alpha_N}, \quad \alpha_N \approx 0.076
    \end{equation}

    \item \textbf{GPT-4}\index{GPT-4} (OpenAI, 2023) : vraisemblablement un mod\`ele Mixture of Experts (MoE)\index{Mixture of Experts} avec des centaines de milliards de param\`etres.

    \item \textbf{Claude}\index{Claude} (Anthropic) : famille de mod\`eles utilisant le RLHF (Reinforcement Learning from Human Feedback)\index{RLHF} et le Constitutional AI.

    \item \textbf{LLaMA}\index{LLaMA} (Meta, 2023--2024) : mod\`eles ouverts utilisant Pre-Norm avec RMSNorm\index{RMSNorm}, RoPE\index{RoPE} (Rotary Position Embeddings), SwiGLU\index{SwiGLU} au lieu de ReLU, et Grouped Query Attention\index{GQA}.

    \item \textbf{Mixture of Experts (MoE)}\index{MoE} : Mixtral (Mistral, 2024) active seulement 2 experts sur 8 par token, permettant un mod\`ele de 47B param\`etres avec un co\^ut d'inf\'erence de $\sim$13B.
\end{itemize}
\end{sota}

\begin{architecturebox}[title=\textbf{\'Evolution des architectures Transformer}]
\begin{center}
\begin{tikzpicture}[
    node distance=1.4cm,
    yearnode/.style={rectangle, draw, fill=blue!10, minimum width=2.5cm, minimum height=0.7cm, align=center, rounded corners=2pt, font=\small},
    arr/.style={-{Stealth}, thick},
]
\node[yearnode] (t0) {Transformer\\(2017)};
\node[yearnode, right=1.5cm of t0] (bert) {BERT\\(2018)};
\node[yearnode, above=0.8cm of bert] (gpt) {GPT / GPT-2\\(2018--2019)};
\node[yearnode, right=1.5cm of bert] (vit) {ViT\\(2020)};
\node[yearnode, right=1.5cm of gpt] (gpt3) {GPT-3\\(2020)};
\node[yearnode, right=1.5cm of vit] (llama) {LLaMA / Mistral\\(2023--2024)};
\node[yearnode, right=1.5cm of gpt3] (gpt4) {GPT-4 / Claude\\(2023--2025)};

\draw[arr] (t0) -- (bert);
\draw[arr] (t0) |- (gpt);
\draw[arr] (bert) -- (vit);
\draw[arr] (gpt) -- (gpt3);
\draw[arr] (gpt3) -- (gpt4);
\draw[arr] (vit) -- (llama);
\draw[arr] (gpt3) -- (llama);
\end{tikzpicture}
\end{center}
\end{architecturebox}

\section{Exercices}
\label{sec:transformers-exercices}

\begin{exercice}[Calcul de la complexit\'e de l'attention]
\label{exo:attention-complexity}
\textbf{Difficult\'e~: 1/3}
Montrez que la complexit\'e temporelle de l'attention scaled dot-product est $O(n^2 d_k)$ o\`u $n$ est la longueur de la s\'equence et $d_k$ la dimension des cl\'es. Pourquoi cela pose-t-il probl\`eme pour les longues s\'equences~? Citez deux approches pour r\'eduire cette complexit\'e.
\end{exercice}

\begin{exercice}[Masque causal]
\label{exo:causal-mask}
\textbf{Difficult\'e~: 1/3}
Impl\'ementez en PyTorch une fonction qui g\'en\`ere un masque causal de taille $n \times n$ et montrez son effet sur la matrice d'attention pour une s\'equence de 5 tokens.
\end{exercice}

\begin{exercice}[Nombre de param\`etres d'un Transformer]
\label{exo:transformer-params}
\textbf{Difficult\'e~: 2/3}
Pour un Transformer encodeur \`a $L$ couches, $h$ t\^etes, dimension $d_{\text{model}}$ et FFN de dimension $d_{ff}$, calculez le nombre total de param\`etres (sans compter les embeddings).

\emph{Indice} : consid\'erez s\'epar\'ement l'attention multi-t\^ete ($4 d_{\text{model}}^2$), le FFN ($2 d_{\text{model}} d_{ff}$), et les LayerNorm ($4 d_{\text{model}}$) par bloc.
\end{exercice}

\begin{exercice}[Impl\'ementation de ViT]
\label{exo:vit-impl}
\textbf{Difficult\'e~: 3/3}
Modifiez le code de la section~\ref{sec:transformer-pytorch} pour cr\'eer un Vision Transformer :
\begin{enumerate}
    \item Impl\'ementez le \texttt{PatchEmbedding} qui d\'ecoupe une image $32 \times 32 \times 3$ en patches de taille $8 \times 8$.
    \item Ajoutez un token \texttt{[CLS]} appris.
    \item Entra\^inez sur CIFAR-10 et rapportez la pr\'ecision.
\end{enumerate}
\end{exercice}

\begin{exercice}[Comparaison BERT vs GPT\@]
\label{exo:bert-vs-gpt}
\textbf{Difficult\'e~: 2/3}
\begin{enumerate}
    \item Expliquez pourquoi BERT ne peut pas \^etre utilis\'e directement pour la g\'en\'eration de texte.
    \item Montrez math\'ematiquement que l'objectif MLM de BERT n'est pas \'equivalent \`a maximiser $\log p(x)$ tandis que l'objectif autoregressif de GPT l'est.
    \item Dans quel sc\'enario pr\'ef\'ererait-on BERT \`a GPT et vice versa~?
\end{enumerate}
\end{exercice}

\begin{exercice}[Analyse des t\^etes d'attention {\normalfont (projet)}]
\label{exo:attention-heads}
\textbf{Difficult\'e~: 3/3}
Entra\^inez un Transformer \`a 4 couches et 8 t\^etes sur une t\^ache de classification de texte. Visualisez les matrices d'attention de chaque t\^ete pour des phrases d'exemple. Identifiez-vous des t\^etes sp\'ecialis\'ees (position, syntaxe, s\'emantique)~? \'Evaluez l'impact de l'\'elagage de certaines t\^etes sur la performance.
\end{exercice}
